\chapter{Metric, Data, and Execution Provenance}

This appendix identifies the 32 registered result fields and records how the final evidence was produced. These fields represent 28 distinct calculations because four fields are reporting aliases. The machine-readable counterpart is \path{ekg/final-frozen-evidence-2026-08-07/final-evidence-manifest.json}. It hashes every accepted result, its input manifest, the common evaluator source snapshot, the comparator summary, and the final radar figure.

\section{Definitive Registered-Result Inventory}
\label{app:core_metric_inventory}

The following tables are generated from the frozen implementation's list of registered result fields and checked against the implementation. The 32 fields represent 28 distinct calculations because four are aliases used by another reporting module. Other JSON fields provide counts, denominators, status values, samples, or diagnostic details rather than additional metrics.

\begingroup\scriptsize\setlength{\tabcolsep}{3pt}\setlength{\LTpre}{0.4em}\setlength{\LTpost}{0.4em}
\begin{longtable}{p{0.06\textwidth}p{0.25\textwidth}p{0.61\textwidth}}
\caption{Registered result-field names and dimensions.}\label{tab:core_metric_inventory_names}\\
\toprule ID & Dimension & Metric \\
\midrule\endfirsthead
\toprule ID & Dimension & Metric \\
\midrule\endhead
\bottomrule\endfoot
M01 & Graph connectivity and structure & Weak/undirected connected components \\
M02 & Graph connectivity and structure & Giant component ratio \\
M03 & Graph connectivity and structure & Average local clustering coefficient \\
M04 & Graph connectivity and structure & Edge connectivity \\
M05 & Graph connectivity and structure & Graph density \\
M06 & Graph connectivity and structure & Average degree \\
M07 & Redundancy and duplication & Exact-label duplicate candidate rate \\
M08 & Redundancy and duplication & Label uniqueness rate \\
M09 & Redundancy and duplication & Fuzzy duplicate candidate pairs \\
M10 & Temporal consistency & Temporal literal validity rate \\
M11 & Temporal consistency & Temporal granularity distribution \\
M12 & Temporal consistency & Temporal coverage \\
M13 & Temporal consistency & Interval validity rate \\
M14 & Temporal consistency & Temporal span \\
M15 & Temporal consistency & Average events per decade in observed span \\
M16 & Minimal event-profile alignment & Label coverage \\
M17 & Minimal event-profile alignment & Date coverage \\
M18 & Minimal event-profile alignment & Minimal event-profile conformance rate \\
M19 & Minimal event-profile alignment & Reported non-allowlisted property count (cap 50) \\
M20 & Completeness & Event-class schema coverage \\
M21 & Completeness & Minimal population completeness \\
M22 & Completeness & Profile label coverage \\
M23 & Completeness & Profile temporal coverage \\
M24 & Completeness & Profile place coverage \\
M25 & Type consistency & Closed-profile explicit type-alignment summary \\
M26 & Entity richness & Average distinct-predicate richness \\
M27 & Entity richness & Sparse event rate \\
M28 & External mapping coverage & owl:sameAs external-link coverage \\
M29 & External mapping coverage & Wikidata link coverage \\
M30 & External mapping coverage & DBpedia link coverage \\
M31 & Predicate usage & Normalised predicate Shannon entropy \\
M32 & Predicate usage & Predicate HHI concentration \\
\end{longtable}

\begin{longtable}{p{0.06\textwidth}p{0.38\textwidth}p{0.48\textwidth}}
\caption{Registered JSON result fields and implementation paths.}\label{tab:core_metric_inventory_paths}\\
\toprule ID & JSON output path & Implementation path \\
\midrule\endfirsthead
\toprule ID & JSON output path & Implementation path \\
\midrule\endhead
\bottomrule\endfoot
M01 & \path{num_components} & \path{ekg_eval_cli.analyzer.GraphAnalyzer.calculate_metrics} \\
M02 & \path{giant_component_ratio} & \path{ekg_eval_cli.analyzer.GraphAnalyzer.calculate_metrics} \\
M03 & \path{avg_clustering} & \path{ekg_eval_cli.analyzer.GraphAnalyzer.calculate_metrics} \\
M04 & \path{edge_connectivity} & \path{ekg_eval_cli.analyzer.GraphAnalyzer.calculate_metrics} \\
M05 & \path{density} & \path{ekg_eval_cli.analyzer.GraphAnalyzer.calculate_metrics} \\
M06 & \path{avg_degree} & \path{ekg_eval_cli.analyzer.GraphAnalyzer.calculate_metrics} \\
M07 & \path{redundancy.duplication_rate} & \path{ekg_eval_cli.redundancy.RedundancyAnalyzer.analyze_redundancy} \\
M08 & \path{redundancy.label_quality.label_uniqueness_rate} & \path{ekg_eval_cli.redundancy.RedundancyAnalyzer.analyze_redundancy} \\
M09 & \path{redundancy.fuzzy_duplicate_pairs} & \path{ekg_eval_cli.redundancy.RedundancyAnalyzer.analyze_redundancy} \\
M10 & \path{temporal.date_format_validation.compliance_rate} & \path{ekg_eval_cli.temporal.TemporalValidator.validate_temporal_consistency / ekg_eval_cli.temporal.TemporalValidator.validate_temporal_semantics} \\
M11 & \path{temporal.temporal_granularity.granularity_percentages} & \path{ekg_eval_cli.temporal.TemporalValidator.validate_temporal_consistency / ekg_eval_cli.temporal.TemporalValidator.validate_temporal_semantics} \\
M12 & \path{temporal.temporal_coverage.temporal_coverage_rate} & \path{ekg_eval_cli.temporal.TemporalValidator.validate_temporal_consistency / ekg_eval_cli.temporal.TemporalValidator.validate_temporal_semantics} \\
M13 & \path{temporal.semantic_validation.consistency_rate} & \path{ekg_eval_cli.temporal.TemporalValidator.validate_temporal_consistency / ekg_eval_cli.temporal.TemporalValidator.validate_temporal_semantics} \\
M14 & \path{temporal.temporal_density.temporal_span_years} & \path{ekg_eval_cli.temporal.TemporalValidator.validate_temporal_consistency / ekg_eval_cli.temporal.TemporalValidator.validate_temporal_semantics} \\
M15 & \path{temporal.temporal_density.avg_events_per_decade} & \path{ekg_eval_cli.temporal.TemporalValidator.validate_temporal_consistency / ekg_eval_cli.temporal.TemporalValidator.validate_temporal_semantics} \\
M16 & \path{schema.label_coverage_rate} & \path{ekg_eval_cli.schema_analyzer.SchemaAnalyzer.analyze_schema_conformance} \\
M17 & \path{schema.date_coverage_rate} & \path{ekg_eval_cli.schema_analyzer.SchemaAnalyzer.analyze_schema_conformance} \\
M18 & \path{schema.schema_conformance_rate} & \path{ekg_eval_cli.schema_analyzer.SchemaAnalyzer.analyze_schema_conformance} \\
M19 & \path{schema.non_standard_properties_count} & \path{ekg_eval_cli.schema_analyzer.SchemaAnalyzer.analyze_schema_conformance} \\
M20 & \path{completeness.schema_coverage_percentage} & \path{ekg_eval_cli.completeness.CompletenessAnalyzer.analyze_completeness / ekg_eval_cli.completeness.CompletenessAnalyzer.analyze_population_completeness} \\
M21 & \path{completeness.population_completeness_percentage} & \path{ekg_eval_cli.completeness.CompletenessAnalyzer.analyze_completeness / ekg_eval_cli.completeness.CompletenessAnalyzer.analyze_population_completeness} \\
M22 & \path{completeness.population_completeness.label_coverage_rate} & \path{ekg_eval_cli.completeness.CompletenessAnalyzer.analyze_completeness / ekg_eval_cli.completeness.CompletenessAnalyzer.analyze_population_completeness} \\
M23 & \path{completeness.population_completeness.temporal_coverage_rate} & \path{ekg_eval_cli.completeness.CompletenessAnalyzer.analyze_completeness / ekg_eval_cli.completeness.CompletenessAnalyzer.analyze_population_completeness} \\
M24 & \path{completeness.population_completeness.location_coverage_rate} & \path{ekg_eval_cli.completeness.CompletenessAnalyzer.analyze_completeness / ekg_eval_cli.completeness.CompletenessAnalyzer.analyze_population_completeness} \\
M25 & \path{type_consistency.overall_type_consistency} & \path{ekg_eval_cli.type_consistency.TypeConsistencyAnalyzer.analyze_type_consistency} \\
M26 & \path{entity_richness.avg_properties_per_event} & \path{ekg_eval_cli.entity_richness.EntityRichnessAnalyzer.analyze_entity_richness} \\
M27 & \path{entity_richness.sparse_entities_percentage} & \path{ekg_eval_cli.entity_richness.EntityRichnessAnalyzer.analyze_entity_richness} \\
M28 & \path{mapping_coverage.external_link_rate} & \path{ekg_eval_cli.mapping_coverage.MappingCoverageAnalyzer.analyze_mapping_coverage} \\
M29 & \path{mapping_coverage.wikidata_coverage} & \path{ekg_eval_cli.mapping_coverage.MappingCoverageAnalyzer.analyze_mapping_coverage} \\
M30 & \path{mapping_coverage.dbpedia_coverage} & \path{ekg_eval_cli.mapping_coverage.MappingCoverageAnalyzer.analyze_mapping_coverage} \\
M31 & \path{predicate_usage.normalized_shannon_entropy} & \path{ekg_eval_cli.predicate_usage.PredicateUsageAnalyzer.analyze_predicate_usage} \\
M32 & \path{predicate_usage.hhi_concentration} & \path{ekg_eval_cli.predicate_usage.PredicateUsageAnalyzer.analyze_predicate_usage} \\
\end{longtable}

\begin{longtable}{p{0.05\textwidth}p{0.37\textwidth}p{0.22\textwidth}p{0.25\textwidth}}
\caption{Exact formula contracts, unavailable-value policies, and provenance classification.}\label{tab:core_metric_inventory_contracts}\\
\toprule ID & Formula / operational definition & Empty or inapplicable case & Source / formalisation class \\
\midrule\endfirsthead
\toprule ID & Formula / operational definition & Empty or inapplicable case & Source / formalisation class \\
\midrule\endhead
\bottomrule\endfoot
M01 & Number of maximal connected components in the undirected projected graph. & No projected nodes: reject the input and emit no structural profile. & graph-theory standard: Standard connected-component definition; NetworkX connected\_components implementation. \\
M02 & |V\_max| / |V|. & No projected nodes: reject the input and emit no structural profile. & adapted from graph theory: Standard giant-component concept normalised by graph order. \\
M03 & (1 / |V|) * sum\_v (2e\_v / (k\_v(k\_v - 1))) on the undirected simple projection. & No projected nodes: reject the input. In large-graph mode, clustering is unavailable and accompanied by an explicit status. & graph-theory standard: Watts-Strogatz clustering coefficient; NetworkX average\_clustering implementation. \\
M04 & Minimum number of edges whose removal disconnects the graph; set to 0 for disconnected graphs. & No projected nodes: reject the input. In large-graph mode, report only conditionally exact cases; otherwise mark the value unavailable. & graph-theory standard: Standard edge-connectivity definition; NetworkX edge\_connectivity implementation. \\
M05 & 2m / (n(n - 1)) for the undirected simple projection. & No projected nodes: reject the input and emit no structural profile. & graph-theory standard: Standard graph-density formula; NetworkX density implementation. \\
M06 & 2m / n for the undirected simple projection. & No projected nodes: reject the input and emit no structural profile. & graph-theory standard: Handshaking lemma; standard graph degree definition. \\
M07 & events in repeated normalised-label groups / events with eligible labels * 100. & No eligible labels: rate is null (N/A); duplicate candidate counts remain 0. & custom operationalisation adapted from conciseness: Linked Data conciseness dimension; thesis-specific exact-label duplicate proxy. \\
M08 & unique normalised sampled labels / sampled labels * 100. & No eligible labels: rate is null (N/A); duplicate candidate counts remain 0. & custom descriptive metric: Motivated by duplicate-detection and label-quality review practice. \\
M09 & Exact count of non-identical sampled normalised-label pairs with RapidFuzz token-sort similarity >= 90\%. & No eligible labels: rate is null (N/A); duplicate candidate counts remain 0. & adapted from record-linkage literature: Fellegi-Sunter record linkage; duplicate-detection literature; RapidFuzz token\_sort\_ratio implementation. \\
M10 & valid ISO-parseable temporal literals / sampled temporal literals * 100. & No applicable temporal values or intervals: rate/value is null (N/A), with an explicit status. & standard-based / literature-derived: XML Schema date lexical forms, OWL-Time temporal representation, syntactic validity in KG quality literature. \\
M11 & count(values at granularity g) / sampled temporal values * 100. & No applicable temporal values or intervals: rate/value is null (N/A), with an explicit status. & adapted descriptive temporal metric: OWL-Time supports multiple temporal positions/granularities; event-KG literature emphasises temporal anchoring. \\
M12 & events with sem:hasBeginTimeStamp or sem:hasEndTimeStamp / direct sem:Event instances * 100. & No applicable temporal values or intervals: rate/value is null (N/A), with an explicit status. & adapted from completeness literature: Profile-relative completeness in KG/Linked Data quality; SEM/EventKG event-time role. \\
M13 & events with parseable start/end and end >= start / checked start-end events * 100. & No applicable temporal values or intervals: rate/value is null (N/A), with an explicit status. & custom logical consistency operationalisation: OWL-Time interval semantics; temporal consistency principles. \\
M14 & max(year) - min(year) over dated events. & No applicable temporal values or intervals: rate/value is null (N/A), with an explicit status. & custom descriptive metric grounded in temporal distribution analysis: OWL-Time temporal positions; event dataset coverage analysis. \\
M15 & sum of distinct-event counts per represented begin-date literal / number of decades from the minimum to maximum observed year, including empty decades. & No applicable temporal values or intervals: rate/value is null (N/A), with an explicit status. & custom descriptive metric: Temporal distribution analysis; thesis-specific event-corpus summary. \\
M16 & events with rdfs:label / total sem:Event instances * 100. & No direct sem:Event instances: profile rates are null (N/A); property count remains descriptive. & adapted from completeness literature: Profile-relative completeness in KG quality literature; rdfs:label as RDF labeling convention. \\
M17 & events with sem:hasBeginTimeStamp or sem:hasEndTimeStamp / total direct sem:Event instances * 100. & No direct sem:Event instances: profile rates are null (N/A); property count remains descriptive. & adapted from completeness and event-modelling literature: KG completeness literature; SEM/EventKG event-time modelling. \\
M18 & events with rdfs:label, (sem:hasBeginTimeStamp or sem:hasEndTimeStamp), and sem:hasPlace / total direct sem:Event instances * 100. & No direct sem:Event instances: profile rates are null (N/A); property count remains descriptive. & custom profile-relative conformance formula: Closed application-profile validation principle; completeness literature under explicit profile requirements. \\
M19 & Count of the up to 50 most frequent distinct event predicates whose namespace is outside the configured allow-list. & No direct sem:Event instances: profile rates are null (N/A); property count remains descriptive. & custom descriptive metric: Vocabulary/schema profiling practice; thesis-specific namespace allow-list. \\
M20 & event classes used under sem:Event subclass closure / sem:Event plus declared direct subclasses of sem:Event * 100. & No direct sem:Event instances: event-profile rates are null (N/A); schema coverage remains schema-relative. & adapted from schema completeness literature: Pipino/Farber-style schema completeness adapted to KG profiles. \\
M21 & direct sem:Event instances with label, begin or end time, and place / all direct sem:Event instances * 100. & No direct sem:Event instances: event-profile rates are null (N/A); schema coverage remains schema-relative. & adapted from completeness literature: Profile-relative completeness; column/property completeness adapted to event resources. \\
M22 & events with rdfs:label / total sem:Event instances * 100. & No direct sem:Event instances: event-profile rates are null (N/A); schema coverage remains schema-relative. & adapted from completeness literature: Profile-relative KG completeness. \\
M23 & events with sem:hasBeginTimeStamp or sem:hasEndTimeStamp / total direct sem:Event instances * 100. & No direct sem:Event instances: event-profile rates are null (N/A); schema coverage remains schema-relative. & adapted from completeness literature: Profile-relative KG completeness and event-time role modelling. \\
M24 & events with sem:hasPlace / total sem:Event instances * 100. & No direct sem:Event instances: event-profile rates are null (N/A); schema coverage remains schema-relative. & adapted from completeness and event-modelling literature: Event role completeness from SEM/event-centric KG modelling. \\
M25 & all explicitly aligned domain and range checks / all applicable declared domain and range checks * 100; null if none apply. & No applicable domain/range evidence: score is null (N/A), not 100\%. & custom closed-profile aggregate over adapted consistency checks: Custom summary informed by RDF/RDFS domain and range semantics and closed-profile validation practice. \\
M26 & mean\_e count(distinct outgoing predicates other than rdf:type attached to direct event e). & No direct sem:Event instances: richness and sparsity are null (N/A). & adapted/custom metric: Adapted from OntoQA attribute richness; event-level descriptive richness is thesis-specific. \\
M27 & events with fewer than 3 distinct outgoing predicates excluding rdf:type / analysed events * 100. & No direct sem:Event instances: richness and sparsity are null (N/A). & custom/new metric: Thesis-specific thresholded sparsity indicator. \\
M28 & events with at least one owl:sameAs link / total events * 100. & No direct sem:Event instances: all mapping rates are null (N/A). & literature-derived / standard vocabulary: Linked Data interlinking quality; OWL owl:sameAs identity predicate. \\
M29 & events with owl:sameAs on a Wikidata HTTP(S) host / total direct sem:Event instances * 100. & No direct sem:Event instances: all mapping rates are null (N/A). & custom source-specific interlinking metric: Adapted from Linked Data interlinking coverage by target source. \\
M30 & events with owl:sameAs on a DBpedia HTTP(S) host (including language subdomains) / total direct sem:Event instances * 100. & No direct sem:Event instances: all mapping rates are null (N/A). & custom source-specific interlinking metric: Adapted from Linked Data interlinking coverage by target source. \\
M31 & H / log2(k), where k is the number of predicates. & No predicates: normalised entropy and HHI are null (N/A). & literature-derived from information theory: Normalised Shannon entropy over categorical distributions. \\
M32 & sum\_i p\_i\textasciicircum{}2. & No predicates: normalised entropy and HHI are null (N/A). & literature-derived from concentration measurement: Herfindahl-Hirschman Index applied to predicate frequency distribution. \\
\end{longtable}
\endgroup


\section{Interpretation of Metric Provenance}

The provenance classes in Tables \ref{tab:core_metric_inventory_names}--\ref{tab:core_metric_inventory_contracts} have precise meanings. \emph{Adopted} means that a standard formula is used without changing its mathematical definition, such as graph density, Shannon entropy, or HHI. \emph{Adapted} means that a known quality principle is applied to the declared direct-event population, such as completeness or interlinking coverage. \emph{Formalised in this thesis} means that prior work motivates the property, but this thesis specifies the exact event population, denominator, threshold, or aggregation rule. The implementation uses the internal labels \emph{literature-informed original}, \emph{custom}, and \emph{project-specific} for this last class. These labels do not claim that general ideas such as completeness, conciseness, or temporal validity were invented here \citep{newman2018,zaveri2016,shannon1948,gini1912,hirschman1964}.

The inventory contains 32 registered result fields, labelled M01--M32 for reference. They represent 28 distinct calculations. M16/M22 repeat label coverage, M12/M17/M23 repeat temporal coverage, and M18/M21 repeat the same label--time--place profile rate. These aliases are retained so each reporting module can display the value it needs, but they are interpreted once and never added together. Other exported fields provide numerators, denominators, sample sizes, status values, or diagnostic details; they are not additional metrics.

\section{Synthetic Dataset Generation}

The final controlled inputs were generated by \path{ekg/create_tiered_synthetic_eventkgs.py}. The generator uses fixed pseudo-random seeds and writes EventKG-style RDF in N-Triples, together with schema-support files. D1 is generated directly as a high-quality profile. D2 and D3 are independently generated from declared profiles rather than produced by repeatedly mutating the same physical files. Their parameters intentionally reduce selected labels, dates, places, descriptions, relations, and external links, while increasing label reuse, temporal coarsening, and type mismatch. This construction supplies known perturbations but also creates the sensitivity/independence limitation discussed in Chapter 5.

Table \ref{tab:appendix_synthetic_generation_parameters} provides the complete profile values used in the final generation run.

\begin{table}[htbp]
\centering
\scriptsize
\caption{Synthetic generation parameters used for the final inputs.}
\label{tab:appendix_synthetic_generation_parameters}
\begin{tabular}{lccc}
\toprule
Parameter & D1 high & D2 mixed & D3 low \\
\midrule
Core direct events & 100 & 120 & 120 \\
Additional text events & 40 & 35 & 25 \\
Entity resources & 320 & 180 & 80 \\
Label probability & 1.00 & 0.75 & 0.45 \\
Date probability & 1.00 & 0.70 & 0.40 \\
Place probability & 0.95 & 0.55 & 0.10 \\
External-link probability & 1.00 & 0.60 & 0.12 \\
Description probability & 0.90 & 0.50 & 0.10 \\
Relation probability & 0.95 & 0.65 & 0.20 \\
Relations/event & 3--5 & 1--3 & 0--1 \\
Unique event-name pool & 100 & 55 & 8 \\
Clustered dates & No & No & Yes \\
Year-only date probability & 0.00 & 0.20 & 0.55 \\
\bottomrule
\end{tabular}
\end{table}

The proposed evaluator and every external comparator consumed the same 16 N-Triples source files per profile. \path{ekg/tool-comparison/corrected-2026-08-07/input-manifest.json} lists each filename, byte size, and SHA-256 hash. The final bundle builder compares these entries with the input manifests embedded in the proposed-framework JSON and fails if any file differs. Database folders, prior results, and benchmark helper files are not included in the input hash.

\section{Repository Reproduction Entry Point}

The evaluation repository is available at \url{https://github.com/tadiwa-aizen/ekg}. Its \path{REPRODUCE.md} file gives the current prerequisites, commands, expected outputs, and disk requirements. From the repository root, the recommended examiner command is:

\begin{verbatim}
.\reproduce.ps1 -Mode Quick
\end{verbatim}

This command creates a separate output folder, runs the tests, regenerates D1--D3, evaluates them, and checks all 32 registered fields against the accepted evidence. \texttt{Verify} checks the stored evidence without recomputation. \texttt{Full} additionally runs ChronoGrapher and OEKG. The report is written under \path{reproduction-output/<timestamp>-<mode>/report/} in Markdown, HTML, and JSON formats.

\section{Frozen Evaluator and Validation Environment}

This section records the exact code and software used to produce the final results. It allows an examiner to confirm that the reported values came from this version of the evaluator and to rerun the same checks without guessing the required dependencies.

The evaluator source snapshot has SHA-256 prefix \FrozenEvaluatorHash. The full hash and every source-file hash are embedded in each result JSON. The runtime was Python 3.12.6 with Click 8.1.7, NetworkX 3.3, RDFLib 7.6.0, RapidFuzz 3.14.5, DuckDB 1.5.5, NumPy 1.26.4, and Numba 0.66.0. Exact Python dependencies are pinned in \path{ekg/ekg-eval-cli/requirements-lock.txt}. Apache Jena and Fuseki 5.6.0 were used, and the Jena archive was checked against Apache's published SHA-512 checksum.

The complete evaluation command is recorded in \path{ekg/run_final_first_party_evaluations.ps1}. It runs D1--D3, the three ChronoGrapher artefacts, and OEKG, then regenerates the radar figure, result macros, metric inventory, and evidence manifest. The \texttt{-SkipOekg} option is only for shorter development checks; the final evidence includes the full OEKG run.

The project is licensed under the MIT licence in \path{ekg/ekg-eval-cli/LICENSE}. The final automated suite is run with:

\begin{verbatim}
cd ekg/ekg-eval-cli
python -m compileall -q ekg_eval_cli tests
python -m pytest -q
\end{verbatim}

The final suite contains 24 passing tests. It covers all 32 registered result-field contracts through hand-computable fixtures and additionally checks projection semantics, deterministic duplicate normalisation, temporal edge cases, unavailable values, direct-event denominators, type applicability, mapping hosts, source hashes, database safety, and cache invalidation.

\section{External Comparator Reproducibility}
\label{app:comparator_reproducibility}

The external comparison uses native outputs and within-tool interpretation. No score is rescaled by the maximum among D1--D3, and no mean is formed across tools. The complete machine-readable record is \path{ekg/tool-comparison/corrected-2026-08-07/comparator-summary.json}; raw logs and outputs are stored beside it.

Table \ref{tab:appendix_comparator_versions} records the exact revisions and local execution adaptations needed to reproduce those native outputs.

\begin{table}[htbp]
\centering
\scriptsize
\caption{Executed comparator revisions, licences, and local adaptations.}
\label{tab:appendix_comparator_versions}
\begin{tabularx}{\textwidth}{>{\raggedright\arraybackslash}p{2.0cm}>{\raggedright\arraybackslash}p{3.1cm}>{\raggedright\arraybackslash}p{1.5cm}X}
\toprule
Tool & Git revision & Licence & Local execution adaptation \\
\midrule
RDFUnit & \texttt{4d98ab1c488f} & Apache-2.0 & None; standalone test generation/execution on local file inputs. \\
KGHeartBeat & \texttt{85683e1fbb20} & MIT & Accepted local endpoints, disabled public-catalogue lookups and local-URI dereferencing, tolerated absent optional LLM imports, and serialised FAIR output. Patch hash \texttt{534fa3de1b56}. \\
Luzzu Framework & \texttt{9a117e7f2c0b} & GPL-3.0 & Local metric/vocabulary configuration and communication-module compatibility changes. Framework patch hash \texttt{f2ebd35dfadb}; metric repository revision \texttt{2124dccc8439} (MIT). \\
SANSA Stack & \texttt{5602abab29e5} & Apache-2.0 & Enabled local Spark execution and isolated metrics so one failure did not abort all outputs. Patch hash \texttt{7c24d5cc65bd}. \\
\bottomrule
\end{tabularx}
\end{table}

The reproducible command entry points are:

\begin{verbatim}
python ekg/tool-comparison/prepare_corrected_inputs.py
python ekg/tool-comparison/run_luzzu_corrected.py
python ekg/tool-comparison/summarize_corrected_comparators.py
\end{verbatim}

RDFUnit, KGHeartBeat, and SANSA were invoked from their pinned local repositories. Their complete command output is retained in \path{ekg/tool-comparison/corrected-2026-08-07/}. The relevant files are \path{rdfunit-*.log}, \path{kgheartbeat-run.log}, and \path{sansa-*.log}. RDFUnit executed 196 generated tests and reported 192 successful test definitions for each profile. Luzzu produced 13 configured metrics for each profile. SANSA returned 14 attempted outputs and failed two metrics on non-literal nodes; those failures are retained rather than replaced. KGHeartBeat's generic overall score was 0.301, 0.304, and 0.307, which does not support the intended quality ordering. These details are why Chapter 4 compares native coverage and sensitivity rather than a common score.

\section{OEKG Cleaning and Complete Event-Layer Run}
\label{app:oekg_run_details}

The source is the OEKG V2.0 Zenodo record (\url{https://zenodo.org/record/4503163}), released on 31 January 2021 and retrieved on 2 August 2026. The working input contains all 13 published event-layer N-Triples files. Seven lines in the cleaned copy of \path{relations_events_literals.nt} contained invalid backslash escapes that prevented standards-compliant parsing. The script \path{ekg/real-oekg-evaluation/clean_oekg_literals.py} repaired only those escapes. The file \path{oekg-provenance-manifest.json} in the same evaluation directory records each changed line and its before/after hash. No triples or semantic assertions were added, and the original extracted files were retained separately. Table \ref{tab:oekg_hashes} records the principal fingerprints; the provenance manifest records every source and repaired-line hash.

\begin{table}[htbp]
\centering
\small
\caption{Principal OEKG source, cleaned-input, and projection fingerprints.}
\label{tab:oekg_hashes}
\begin{tabularx}{\textwidth}{>{\raggedright\arraybackslash}p{3.2cm}X}
\toprule
Artefact & SHA-256 \\
\midrule
Source archive & \texttt{392d6eeb69d074130166fb626d4db7279}\newline\texttt{c2b16b45be50213480de1864ce4aa4a} \\
Cleaned input aggregate & \texttt{fbf2c85305d653916d8c181dba15e262}\newline\texttt{8c4d0ce2ce4aa6cfc3d0b8b8ff86b3ab} \\
Projection implementation & \texttt{c2d0a6fccb1ef88c85b02f1e97ad2149}\newline\texttt{c0279cb7f9f891c1f84c077f81d59c9b} \\
\bottomrule
\end{tabularx}
\end{table}

The TDB2 manifest records 93,474,126 triples and 954,554 direct \texttt{sem:Event} resources. The structural projection manifest records the input hash, projection contract, and projection-code fingerprint shown in Table \ref{tab:oekg_hashes}. Cache reuse is rejected if the input or projection code differs.

The final command shape was:

\begin{verbatim}
python -m ekg_eval_cli.cli oekg-event-layer-clean --verbose
  --large-graph-mode --output-dir final-frozen-evidence-2026-08-07/oekg
  --large-graph-work-dir large-graph-work-oekg
  --duckdb-memory-limit 8GB
  --duckdb-temp-dir large-graph-work-oekg/duckdb-temp
  --jena-home apache-jena-5.6.0-full
  --fuseki-home apache-jena-fuseki-5.6.0 --port 3030
\end{verbatim}

The available machine had approximately 20GB RAM and a 200GB working-disk constraint. The structural path streams the projection, stores edges and component state on disk, and does not materialise the whole graph in NetworkX. Average clustering remains unavailable in this mode. The final result file and its SHA-256 are recorded in the final evidence manifest rather than copied manually into this appendix.

\section{ChronoGrapher Artefact Preparation and Runs}

The public Turtle files were obtained from the ChronoGrapher graph-search repository and validated with Apache Jena RIOT before conversion to UTF-8 N-Triples. An initial UTF-16 conversion produced by PowerShell redirection was discarded. The files are stored in \path{ekg/chronographer-evaluation/raw/}. Three contain direct \texttt{sem:Event} assertions and were evaluated: \path{eventkg_ng.ttl}, \path{search_ng.ttl}, and \path{generation_ng.ttl}. The 11.6MB \path{frame_ng.ttl} file was excluded because it has no direct \texttt{sem:Event} population, so the event-profile denominator would be empty.

Legacy TDB caches without input manifests were retained separately before fresh TDB2 databases were built from the unchanged N-Triples files using the verified complete Jena distribution. The standard evaluator command was then run once per folder with separate Fuseki ports. Each final JSON contains the source-file SHA-256, input aggregate hash, common evaluator source hash, runtime versions, parameters, and the same 32 registered result fields.
